\chapter{Methodology}
\section{Conceptual Framework}
The study has two mutually connected, but experimentally distinct, phases. Phase I develops and evaluates a controlled text-classification artifact without requiring sensitive resident complaints. Phase II describes how that artifact would support a human-governed barangay reporting system. The separation is essential: achieving a test-set score completes a model experiment, whereas operating a public-facing reporting system requires application, governance and safety verification. Figure~\ref{fig:overall} presents these phases with an explicit release gate and a feedback loop that returns reviewed corrections to \emph{offline} retraining instead of changing a live classifier automatically.
\begin{figure}[H]\centering\includegraphics[width=0.96\textwidth,height=0.85\textheight,keepaspectratio]{figures/conceptual_framework.pdf}\caption{Expanded two-phase conceptual framework separating research preparation, model selection, provisional inference, human review and controlled feedback. Arrows distinguish the measured research workflow from the proposed operational workflow; administrator feedback is reviewed offline rather than automatically deployed.}\label{fig:overall}\end{figure}
\subsection{Phase I: Evidence-Based Corpus Development}
The offline phase begins with a source register: national legislation provides examples of community responsibility and safeguards, local reporting guidance describes information useful for incident review, linguistic studies inform complaint speech acts and code-switching, and synthetic-data research informs variation and evaluation controls. These sources \emph{inform} the scenario design. They do not certify each generated sentence as authentic or prove that the proportions are representative. A 118-item researcher-defined taxonomy supplies 60 Valid issue topics, 28 Suspicious clarification reasons and 30 Spam types. Each item receives an explicit scope definition, an intended three-class label, and references to the corresponding evidence type.

The generation stage distinguishes \emph{scenario meaning} from \emph{surface form}. A hypothetical clogged-drain concern can be written as an informal question or a detailed narrative without changing its class. A confusing report may instead refer only to ``the thing I reported before'' with no recoverable subject. A promotional solicitation lacks a community grievance altogether. Generator logic chooses a controlled language group, communication form, narrative length and contextual cues, then constructs the sentence. Every output carries source traceability, seed, topic/review/spam identifier and a scenario family. These metadata fields are excluded from the text-only classifier. Figure~\ref{fig:dataset} details this process.
\begin{figure}[H]\centering\includegraphics[width=0.94\textwidth,height=0.82\textheight,keepaspectratio]{figures/dataset_pipeline.pdf}\caption{Scenario-first data construction and audit pipeline. The label is based on intended communicative meaning, not on incidental length, grammar, sentiment, or language.}\label{fig:dataset}\end{figure}
\subsection{Phase II: Proposed Public Reporting Service}
The proposed operational system separates resident input, inference, administrative review, case management, and controlled model maintenance. A resident-facing form collects a required Subject and Complaint Details narrative; incident date/time and Purok/Sitio or an identifying landmark are supporting contextual fields when known. Geographic maps, GPS capture and photographic uploads are outside the selected core thesis scope. The application concept uses an authenticated resident interface and an authorized staff interface; the suggested software stack is Next.js/TypeScript, FastAPI and PostgreSQL, but implementation claims require working source code and testing. No completed user-interface or government-office deployment evaluation is included in the supplied experimental evidence.

The class-prediction output is separate from the administrative disposition. Predicted Valid goes to the accepted-case workflow under the selected routing rule; predicted Suspicious or Spam goes to administrator review. Acceptance, clarification, and rejection are decisions about the submission, not automatic replacements for the text label. Proposed controls include resident status tracking, staff action logs, deadline monitoring, referrals, verification before resolution, and a reconsideration mechanism. Detailed proposed figures and the experimental--operational boundary are presented later in this chapter.

\input{chapters/data_description_and_eda.tex}

\section{Research Design and Research Questions}
The study uses applied, quantitative experimental design for its completed machine-learning work and system-design methodology for the proposed application. The primary output is an auditable three-class text classifier. The dependent variables are classification accuracy, precision, recall, Macro-F1 and subgroup errors; independent experimental choices include input language assignment, category family, textual length, message construction pattern, TF-IDF feature ranges and SVM regularization. No causal effect of human corrections is claimed until a frozen pre-update/post-update evaluation is performed.

\section{Corpus Construction and Provenance}
\subsection{Fixed Research Taxonomy and Label Boundaries}
The 118 categories are a scenario inventory and audit framework, not 118 machine-learning output classes. A Valid record conveys an understandable community-related concern or safety report, even if the barangay must refer it to another agency. Suspicious indicates essential context is missing, self-contradictory, or unrecoverably ambiguous \emph{from the input text}; it is not an accusation of dishonesty. Spam communicates no interpretable barangay-reporting purpose. Missing photographs, the absence of a name, nonstandard grammar, or anger alone do not warrant a Suspicious label. Reports \emph{about} scam incidents differ from scam solicitations \cite{ra7160,ra9262,ibanez}.
\begin{table}[H]\centering\caption{Three-class operational decision rubric}\small\begin{tabularx}{\textwidth}{p{0.21\textwidth}X}\toprule
\textbf{Class} & \textbf{Positive criterion and crucial exception}\\\midrule
Valid & Interpretable event, request or community concern. Short, emotional, anonymous and code-switched complaints may all be Valid.\\
Suspicious & Report-like wording but no recoverable issue, unresolved contradiction, or essential missing reference. It is a \emph{clarification state}, not a finding of falsity.\\
Spam & Promotional, impersonation or irrelevant submissions lacking citizen reporting intent; also nonsensical or classifier-directed text. A citizen reporting spam or scam activity can be Valid.\\\bottomrule
\end{tabularx}\end{table}

\subsection{Source-to-Category Traceability}
The evidence matrix bundled with the Overleaf project maps all 118 category IDs to laws, municipal complaint procedures, literature or researcher-defined rules. The mapping is deliberately qualified: statutes justify topics and jurisdictional context; linguistic papers justify forms of expression; and operational definitions justify unreported or undecidable content. An entry is not considered evidence that a specific synthetically generated narrative occurred in reality. Specialized domains involving violence, children, fraud, utilities and public danger must use appropriate referral practices. An appendix lists the category register used in generation.

\subsection{Generation Controls and Dataset Composition}
The released study corpus was produced with fixed random seed \texttt{20261026}. It contains exactly 24,000 unique text records: 8,000 per classification, and within each class 2,000 per generator-assigned language group (Bisaya, Filipino, English, Mixed). The generator uses researcher-authored source-aware scenario descriptions, reason-specific ambiguous seed messages and non-complaint examples, combined with structured variations in question forms, greetings, follow-up expressions, requested actions, partial location/time cues, emotional language and length bands. Every row is identified as synthetic and unreviewed; its assigned language is not independently verified. The equal distributions are experimentally controlled and cannot estimate true complaint prevalence.
\begin{table}[H]\centering\caption{Corpus structure and recorded partitions}\begin{tabularx}{\textwidth}{p{.43\textwidth}rX}\toprule
\textbf{Dimension}&\textbf{Count}&\textbf{Role}\\\midrule
Valid / Suspicious / Spam&8,000 each&Target labels\\
Bisaya / Filipino / English / Mixed&6,000 each&Generator allocation\\
Valid topics / Suspicious reasons / Spam types&60 / 28 / 30&Scenario metadata\\
Training&19,213&Fit features and select candidates\\
Validation&2,399&Model selection diagnostics\\
Synthetic test&2,388&Final within-generator estimate\\\bottomrule\end{tabularx}\end{table}

\subsection{Data Dictionary and Feature-Exclusion Boundary}
The predictive feature is \texttt{complaint\_text}; the supervised target is \texttt{classification}. Fields such as \path{category_id}, \path{source_reference_ids}, \path{scenario_family_id}, \path{template_family_id}, \path{language_group}, \texttt{split}, and \texttt{label\_rule} exist only for quality review and experimental controls. Feeding them to the SVM would directly expose the desired label or generation family and invalidate the intended text-only evaluation.

\section{Data Quality Assurance and Threats to Validity}
The generator removed exact-duplicate candidate outputs and verified nonempty narratives, target-class counts and category coverage. These technical checks do not establish natural idiom, correct meaning, grammatical coherence or independent annotation agreement. Accordingly the dataset retains flags marking \texttt{language\_verified=No} and \texttt{human\_label\_reviewed=No}. The remaining required manual assessment is an adjudicated stratified sample across 118 categories, three classes, four assigned languages, length bands and borderline pairs. At least two independent reviewers should code labels while blinded to generation metadata, adjudicate disagreements, and report agreement alongside confusion rates; a fluent Cebuano/Filipino review is essential before deployment claims. These are planned review tasks, not completed results.

\section{Experimental Setup and Models}
\subsection{Text Representation and Linear Models}
TF-IDF transforms the complaint narratives into a sparse feature space. Multinomial Naive Bayes serves as the baseline estimator using word unigrams/bigrams; the primary model is \texttt{LinearSVC} with regularization parameter $C$. Both operate solely on complaint narrative text, but their recorded TF--IDF settings are different: the Naive Bayes baseline uses word unigrams/bigrams whereas the chosen SVM uses word unigrams. Thus the existing comparison contrasts \emph{complete pipelines} rather than isolating estimator choice with a strictly identical feature matrix. A matched-representation sensitivity comparison can be conducted on the same training/validation partitions in the next documented run; its results must not be attributed to the previous experiment. A \texttt{Pipeline} fits the vectorizer within each training fold to avoid using validation or test vocabulary in training estimators \cite{sklearn}. The recorded selected SVM uses word unigrams; an isolated exploratory character 3--5-gram model was evaluated later as an ablation. The combined word-unigram/bigram plus character-within-word 3--5-gram representation is a separately evaluated development candidate, not the feature set producing the historical reference SVM results. Similarly, the reported baseline does not supply fitted calibrated class probabilities. A separate combined-feature sigmoid-calibrated candidate was subsequently evaluated on existing development validation only, as reported in Chapter~4. Independent probability reliability evidence remains required before operational probability-based decisions.
\subsection{Partitioning and Cross-Validation}
The study dataset used scenario-family partitioning before feature fitting: 19,213 train, 2,399 validation, 2,388 test. The 18-config SVM grid combined $C\in\{0.1,1,10\}$, TF-IDF word n-gram range in $\{(1,1),(1,2)\}$ and minimum document frequency in $\{2,5,10\}$. Selection used three-fold \texttt{StratifiedGroupKFold} by scenario-family ID with mean Macro-F1. This avoids direct reuse of that ID across folds but does not eliminate shared seed sentences across different family IDs. The best recorded configuration was $C=0.1$, unigrams, \texttt{min\_df=2}, with \texttt{max\_df=0.95} and sublinear term frequency fixed. The observed tie among all 18 grid configurations does \emph{not} demonstrate that tuning has no effect in genuine reporting language.
\begin{figure}[H]\centering\includegraphics[width=.96\textwidth,height=.80\textheight,keepaspectratio]{figures/m2_unified_evaluation.pdf}\caption{Unified experiment allocation and evaluation protocol. Only the 19,213 training examples participate in model selection and five-fold grouped stability checks; the fixed fitted models are evaluated on 2,399 validation and the previously recorded 2,388 synthetic test examples.}\label{fig:evaluation}\end{figure}

\subsection{Evaluation and Learning Diagnostics}
Accuracy reports the fraction of correct labels, while per-class precision and recall expose specific confusions; Macro-F1 gives equal weight to three classes. Further diagnostics include confusion matrices, class and language subgroup metrics, text-length-only performance, nearest-training-example similarity, and a separately drafted 96-message challenge. Learning curves depict accuracy or Macro-F1 as a function of training-set size; these are \emph{not} loss curves over epochs. The final \texttt{LinearSVC} fit does not expose a documented epoch-loss history. Consequently, a separate \texttt{SGDClassifier(loss=hinge)} run recorded epoch-wise training and validation hinge loss. This auxiliary trace supports an optimization discussion but must not be mislabeled as the selected SVM model's own training trace.


\section{Experimental Setup, Reproducibility and Evaluation}
\label{sec:m2}
This section reconciles the Milestone~2 study plan with the experiments actually completed for the current 24,000-narrative synthetic corpus. The preliminary Milestone~2 document proposed a stratified 80/20 split, five-fold tuning, and a larger TF--IDF parameter grid; those were \emph{planned settings}, not the final run. The recorded experiment instead used scenario-family-grouped partitioning (19,213 training, 2,399 validation, 2,388 test), three-fold grouped SVM grid search over 18 configurations, and a separate five-fold grouped stability analysis. The distinction matters for transparent reproducibility.

\subsection{Hardware and Software Environment}
The hardware below was reported in the provided Milestone~2 document as the research development workstation; it is not independently benchmarked in this manuscript. TF--IDF extraction, grouped cross-validation and scikit-learn's \texttt{LinearSVC} use a CPU-oriented implementation; the listed dedicated GPU is not evidence of GPU-accelerated SVM training.
\begin{table}[H]\centering\caption{Documented research workstation and experimental software}\label{tab:m2env}\small\begin{tabularx}{\textwidth}{p{.27\textwidth}X}\toprule
Component & Reported specification\\\midrule
CPU & Intel Core i7-13700HX (13th generation)\\
GPU & NVIDIA GeForce RTX 4060 Laptop GPU; Intel UHD Graphics integrated adapter\\
Memory & 16 GB installed (two 8 GB Samsung DDR5-5600 modules; 15.75 GB usable)\\
Operating system & Windows 11 Home, version 25H2, build 26200.9457\\
Python & 3.13.5\\
Libraries & pandas 2.2.3; NumPy 2.3.5; scikit-learn 1.8.0; SciPy 1.17.0; joblib 1.5.3\\\bottomrule
\end{tabularx}\end{table}
The Python and scikit-learn versions are also recorded in \path{EXPERIMENT_MANIFEST.json}. Other version assertions come from the submitted Milestone~2 environment record and should be checked against a captured environment export if rerunning on a different computer.

\subsection{Version Control and Reproducibility}
The earlier Milestone~2 submission identifies a GitHub repository at \url{https://github.com/rencze/nlp-barangay-complaint-classification-hitl}. The earlier manuscript recorded an HTTP 404 observation; remote visibility was not reverified during this local reproduction. The repository's visibility, owner permissions and current commit SHA must be verified before submission, and a working read-only link supplied to examiners.

This Overleaf bundle contains \path{README.md}, \path{requirements.txt}, \path{EXPERIMENT_MANIFEST.json}, \path{data/ALL_SYNTHETIC_V2.csv}, \path{reproducibility_scripts/}, and \path{validation_suite/}. To reproduce a \emph{minimal} training example without conflating it with the reported 24,000-row research results, create a Python 3.13 environment, install \texttt{pip install -r requirements.txt}, then run \texttt{python} \path{examples/minimal_pipeline.py}. To reproduce the full research suite, first inspect \path{README.md}, the dataset audit record, the saved split metadata, the bundled experiment scripts and the recorded seed \texttt{20261026}. The source's earlier seed \texttt{42} was only proposed in Milestone~2 and is not the seed of these recorded runs. Model training code, feature fitting, partition membership, candidate settings, dependency versions and software revisions should be preserved together. Do not claim that \texttt{src/train.py} exists in this Overleaf archive: the milestone's separate repository structure was not included here.

\subsection{Baseline and Measured Model}
The standard baseline is TF--IDF plus \texttt{MultinomialNB(alpha=1.0)}, a conventional, relatively inexpensive model for sparse text features. The comparison model is word-level TF--IDF plus a linear support vector classifier, \texttt{LinearSVC}. Both are evaluated with identical training, validation and synthetic-test partitions; their target labels are Valid, Suspicious and Spam. The historical scores compare two predeclared text-processing pipelines with the same partition memberships, not estimators operating on exactly equal word features. A fair estimator-only follow-up can use the \emph{same} word n-gram representation in both models, fitted on the existing training partition and assessed through training-group folds and the existing development validation partition. Such a follow-up is not part of the recorded results and must be labeled independently. The separately proposed concatenated word-plus-character representation is \emph{not} the model whose reported test scores appear in Chapter~4.

\subsection{Evaluation Metrics and Equations}
Let $TP_c$, $FP_c$, $FN_c$ and $TN_c$ denote one-versus-rest confusion counts for class $c$, where $c\in\{\mathrm{Valid,Suspicious,Spam}\}$, and let $N$ denote the number of evaluated examples. The prespecified primary comparative metric is \emph{Macro-F1}, supplemented by accuracy, per-class metrics, weighted-F1 and confusion matrices:\vspace{-4pt}
\begin{align}
\mathrm{Accuracy} &= \frac{\sum_{c} TP_c}{N},\\
\mathrm{Precision}_c &= \frac{TP_c}{TP_c+FP_c},&
\mathrm{Recall}_c &= \frac{TP_c}{TP_c+FN_c},\\
F_{1,c} &= \frac{2\,\mathrm{Precision}_c\,\mathrm{Recall}_c}{\mathrm{Precision}_c+\mathrm{Recall}_c},\\
F_{1,\mathrm{macro}} &= \frac{1}{3}\sum_c F_{1,c},&
F_{1,\mathrm{weighted}} &= \frac{1}{N}\sum_c n_cF_{1,c}.
\end{align}
Here $n_c$ is the number of true records in class $c$; zero-denominator class scores are explicitly set to zero in reporting. Because the current data are deliberately class-balanced, high accuracy is still not proof of generalization. Macro-F1 gives each text label equal influence, while per-class precision and recall detect errors concealed by a single aggregate score. Particular attention is assigned to \emph{Valid predicted as Spam}, potential false rejection risk, and the class-wise confusion matrix. The operational prevalence of these classes is not estimated by the synthetic class allocation.

\subsection{Actual Hyperparameter Search and Selection}
The recorded SVM search used \texttt{GridSearchCV} with three-fold, scenario-family-grouped cross-validation on the 19,213 training records. It examined 18 combinations (three $C$ values by two word-ngram choices by three minimum-document-frequency values); all learned TF--IDF statistics were fitted within training folds. The later, separate five-fold \texttt{StratifiedGroupKFold} experiment retested the selected settings for stability; it did not perform the original search again.
\begin{table}[H]\centering\caption{Actual fixed and varied settings (superseding the preliminary Milestone 2 grid)}\label{tab:m2grid}\small\begin{tabularx}{\textwidth}{p{.39\textwidth}X}\toprule
Setting & Recorded value(s)\\\midrule
\texttt{LinearSVC.C} (variable) & $\{0.1,1,10\}$\\
\texttt{TfidfVectorizer.ngram\_range} (variable) & $\{(1,1),(1,2)\}$ (word analyzer)\\
\texttt{TfidfVectorizer.min\_df} (variable) & $\{2,5,10\}$\\
\texttt{max\_df} / \texttt{sublinear\_tf} (fixed) & $0.95$ / enabled\\
Primary tuning metric & Grouped three-fold mean Macro-F1\\
Selection from tied candidate scores & $C=0.1$, word unigrams, \texttt{min\_df}=2 (deterministic tie handling)\\
Follow-up stability analysis & Separate five-fold grouped CV on training partition\\
Random seed & \texttt{20261026}\\\bottomrule
\end{tabularx}\end{table}
The grouped grid is selected by mean Macro-F1 with deterministic tie handling; observed scores and their interpretation are reported in Chapter~4. The reserved evaluation records were not used to fit TF--IDF or select this grid. The unified allocation and validation protocol is summarized in Figure~\ref{fig:evaluation}.

\subsection{Statistical Model Comparison}
Because the SVM and Naive Bayes predict the \emph{same} held-out narratives, their correct-versus-incorrect outcomes form paired binary observations. McNemar's test is therefore appropriate for comparing error occurrence between \emph{two} classifiers on a common test partition \cite{mcnemar_nist}. Let $b$ be the count correctly classified by the baseline but wrongly by SVM, and $c$ the count wrongly classified by the baseline but correctly by SVM. The continuity-corrected large-sample statistic and the small-discordance exact alternative are:
\begin{align}
\chi^2_{\mathrm{McN}} &= \frac{\max(0,|b-c|-1)^2}{b+c},\quad b+c>0,\\
X&\sim\mathrm{Binomial}(b+c,\tfrac12),\quad
p_{\mathrm{exact}}=\min\!\left\{1,2\Pr\big[X\leq\min(b,c)\big]\right\}.
\end{align}
The hypothesis $H_0$ is equal paired error probabilities, tested at $\alpha=0.05$. The exact binomial form is preferred for the present experiment's small number of discordant pairs. Actual paired predictions, contingency counts and interpretations are reported in Chapter~4. The repeated historical test is not a new independent test. A Friedman test would concern a different design involving more classifiers and repeated comparable datasets; it is not justified by the current two-model, one-test-set comparison.

\subsection{Model Evaluation Sequence}
The actual reproducible sequence is (1) preserve the 24,000-row synthetic source, codebook and scenario-group IDs; (2) group-partition into 19,213/2,399/2,388; (3) fit the baseline and perform three-fold grouped SVM grid search on training only; (4) conduct a separate five-fold grouped stability analysis on the training portion; (5) evaluate the selected fixed models on 2,399 validation records and the previously evaluated 2,388 synthetic test records; (6) report class-wise and aggregate scores, confusion matrices and exact paired-error comparison; and (7) document scripts, seed, versions, selected parameters, generator limitations and the limitation that performance outside this synthetic corpus is not measured. Chapter~4 contains \emph{observed results}; combined features and calibration are separately evaluated development candidates, while operational HITL and model retraining remain proposed system capabilities.

\section{Preserved Reproduction Protocol and Evidence Boundary}
\label{sec:prerun}
For the completed rerun, the researchers preserved the supplied \texttt{ALL\_SYNTHETIC\_V2.csv} and its recorded partition memberships. \textbf{The current repeatability exercise uses the same 24,000-record synthetic corpus; no additional dataset or relabeling is part of this reproduction.} The set comprises 19,213 rows for fitting and grouped cross-validation, 2,399 rows for fixed development validation, and 2,388 rows for the historically evaluated synthetic test. Only the training partition is used for estimator and TF--IDF fitting. A validation score is not a training score, and a repeated score from the same test rows is not new external evidence.

The GitHub artifact should record a Git commit identifier, the SHA-256 of the 24,000-row source file, Python and library versions, experiment seeds and exact commands. The dataset partitions are exported from the \texttt{split} column without generating or resampling records, and a preflight script checks disjoint scenario-family IDs and counts before model fitting. The two principal reproducible commands are \path{python reproducibility_scripts/evaluate_models.py} for the original model comparison and \path{python validation_suite/run_grouped_validation.py} for the separate five-fold training-only stability audit. The repository also includes a small smoke test, which is not evidence of full-study accuracy. Since the full scripts may overwrite recorded outputs, the original result files must be archived and versioned before a rerun. Any discrepancy is logged, not silently rewritten into the thesis.

The first completed task was \emph{reproduction}: confirm that the recorded preprocessing, SVM, Naive Bayes and cross-validation scripts execute on unchanged splits. If a second, explicitly labeled \emph{sensitivity experiment} compares a common TF--IDF representation or the proposed word--character combination, it must use training-group folds and the existing development validation partition for configuration decisions. It must not use the previously viewed synthetic test as an unseen selection set. Report the new settings and results separately instead of replacing the historical scores. A subsequent combined-feature/calibration experiment was executed on development validation only; operational HITL, feedback retraining and live case management remain untested regardless of their appearance in design diagrams.

\subsection{Documentation and Reproduction Checklist}
\begin{table}[H]\centering\caption{Reproducibility artifacts and what each establishes}\label{tab:preruncheck}\small
\begin{tabularx}{\textwidth}{p{.34\textwidth}X}\toprule
Artifact & Required evidence \\\midrule
Source CSV and split files & Source hash; partition totals; label counts; no scenario-family overlap across partitions. \\
Baseline and SVM scripts & Refit TF--IDF within each training fold; record selected parameters and seeds. \\
Validation metrics & Accuracy, class-wise precision, recall, F1, Macro-F1, weighted-F1 and confusion matrices on 2,399 records. \\
Synthetic test archive & Historical scores on 2,388 records, explicitly marked as already evaluated; preserve paired predictions if available. \\
Five-fold stability audit & Training-only grouped folds; group overlap report and fold-wise metrics. \\
README, environment and commit & Installation instructions, exact commands, dependency versions and source-code revision. \\\bottomrule
\end{tabularx}\end{table}

\section{Required Next Validation Protocol}
As \emph{future validation recommendations}, rather than conditions claimed to have been completed in the current single-corpus rerun, robust deployment evaluation would include: (a) independent blinded human label agreement; (b) broader unseen-scenario / root-seed holdout independent of template families; (c) deliberately constructed minimal-pair tests of ambiguous versus clear narratives; (d) subgroup performance on Bisaya, Filipino, English and mixed inputs; (e) safety-sensitive false-negative review; (f) probability/margin reliability and an abstention or ``review required'' policy; and (g) full-stack end-to-end tests for accounts, audit logs, urgency escalation, status transitions, feedback consent, versioning, latency and rollback. These items are not evidence generated by the current experiment. The immediate documented rerun uses only the existing synthetic corpus and its stored train/validation/test membership; any subsequent external assessment, if undertaken, must be reported as a separate study.

\section{System Architecture and Operational Boundary}
The proposed system distinguishes a researcher-evaluated word-level TF--IDF/SVM baseline from a planned operational architecture that adds privacy-aware input preparation, word-plus-character features, calibrated probabilities, class-based HITL review, accepted-case management and controlled retraining. The measured model scores are not evidence that the additional production functions are already implemented or validated. Section~\ref{sec:proposed-operational-workflows} defines these interfaces, states and non-automatic feedback processes. System-level monitoring and administrator accountability are application requirements to be evaluated separately from classifier accuracy.

\section{Proposed Integrated Workflows and Administrative Controls}
\label{sec:proposed-operational-workflows}
\textbf{Implementation boundary.} Figures in this section specify the \emph{target application design}. They are not additional completed experiments. The reproducible results in this manuscript concern the existing synthetic-data word-level TF--IDF and Linear SVM pipeline. The combined word/character representation and grouped sigmoid calibration were separately tested on existing synthetic development validation. Live probability reliability, class-dependent administrative review, case management and feedback-driven retraining still require independent implementation and measurement before production-readiness claims. Each diagram is placed immediately after the explanatory discussion of its named process; figure numbers are assigned by LaTeX captions rather than embedded in the graphics.

\subsection{Complaint Intake and Input Boundary}
The proposed resident-side form requires a \emph{Subject} and a \emph{Complaint Narrative}. If known, a date/time of occurrence and an approximate Purok/Sitio or landmark can be supplied as supporting context, without compulsory GPS or an interactive map. The narrative, not the user's name, phone number, assigned language tag or other metadata, is the primary input to the three-class SVM. The subject can be retained for case identification, but adding it to the classifier requires a separately documented experiment to preserve comparison with the recorded text-only baseline. The application should provide a confirmation and tracking reference and maintain authenticated access and lawful handling of resident information. The screenshots of Davao City Reports serve as a \emph{local interface precedent}, not evidence that its implementation has been replicated \cite{dcr,ra10173}.

The core scope excludes photo uploads, image classification, ensembles that combine images and text, automated determination of physical truth and mandatory GPS coordinates. Where urgent safety reports arise, an independent referral/escalation safeguard is recommended rather than allowing an ordinary three-class text prediction to prevent timely review. Such referral logic remains a system requirement, not an experimentally demonstrated function.

\subsection{NLP Preprocessing: Meaning-Preserving Candidate}
Figure~\ref{fig:op-nlp} specifies a proposed conservative NLP routine. A secure original can be retained for authorized processing while the model operates on a minimized copy with direct identifiers masked when appropriate. Unicode NFC, whitespace normalization, lowercasing and limited cleanup are intended to standardize representation without discarding \emph{hindi}, \emph{wala}, \emph{no}, \emph{not}, relevant number expressions or incident descriptions. Generic stopword removal, stemming and lemmatization are excluded from this proposed routine to avoid meaning-altering changes to mixed-language complaints. The current reproducibility scripts use scikit-learn TF--IDF's text processing; the separately depicted privacy masking and normalization stages are not asserted to be fully implemented in those recorded experiments. Any masking accuracy must be assessed because over-masking can remove useful context and under-masking can expose personal information \cite{ra10173,sklearnpitfalls}.
\begin{figure}[H]\centering\includegraphics[width=.94\textwidth,height=.69\textheight,keepaspectratio]{figures/nlp_proposed.pdf}\caption{Proposed Natural Language Processing (NLP) preprocessing workflow for complaint narratives. Privacy-aware masking and normalization are target features, not measured preprocessing steps in the reported synthetic-data experiments.}\label{fig:op-nlp}\end{figure}

\subsection{TF--IDF Feature Extraction and Model-Level Consistency}
The recorded main SVM results used word-unigram TF--IDF, with a separately examined character-only ablation. For subsequent comparison, Figure~\ref{fig:op-tfidf} proposes word unigrams and bigrams \emph{together with} character-within-word 3--5-grams. The resulting two sparse feature matrices are concatenated before classification. This combination is a \emph{separately evaluated development candidate}, not the verified configuration yielding the previously reported perfect within-generator score. In both configurations, vocabulary and inverse-document-frequency weights must be fitted within training partitions, never independently refitted on each new complaint. The deployed vectorizer and classifier should be versioned together \cite{sklearntfidf,sklearnpitfalls}.
\begin{figure}[H]\centering\includegraphics[width=.96\textwidth,height=.68\textheight,keepaspectratio]{figures/tfidf_proposed.pdf}\caption{Proposed word-level and character-level TF--IDF feature extraction workflow; the concatenated representation remains to be evaluated against the recorded word-only baseline.}\label{fig:op-tfidf}\end{figure}

\subsection{SVM Training and Calibrated Prediction}
Figure~\ref{fig:op-svm} separates model fitting from inference. An offline Linear SVM is trained using authorized labeled data, with hyperparameter selection performed through training-only cross-validation. Before model probabilities can be described as \emph{calibrated}, an appropriate calibrator must be fitted using independent calibration folds or held-out training material, then evaluated for reliability. During inference, a new complaint is transformed by the saved compatible vectorizer and classified by the saved model; it is not used to fit the model. The three displayed class-probability quantities remain proposed operational outputs; the reference SVM is uncalibrated, while a separate candidate supplies synthetic development probability diagnostics in Chapter~4 \cite{sklearncalib,sklearnpitfalls}.
\begin{figure}[H]\centering\includegraphics[width=.96\textwidth,height=.72\textheight,keepaspectratio]{figures/svm_proposed.pdf}\caption{Proposed support vector machine (SVM) model-development and new-complaint prediction workflow, with a separate probability-calibration requirement.}\label{fig:op-svm}\end{figure}

\subsection{Human-in-the-Loop Verification and Administrative Disposition}
The initial proposed routing rule is \emph{class-based}: predictions of Valid go directly to accepted-case handling, while predicted Suspicious or Spam are referred to the administrator for Human-in-the-Loop (HITL) review. Calibrated scores, once validated, may help \emph{prioritize} cases within the review queue; no numerical cutoff is fixed by this study. Under this design, an error that predicts Valid will not enter routine HITL, so a separately selected random audit of predicted Valid cases is necessary to estimate unobserved misclassifications. This is a proposed mitigation of the selective-label problem, not an implemented audit \cite{lakkaraju2017,madras2018}.

The administrator first examines the particular original text version and may confirm or correct its \emph{text label} to Valid, Suspicious or Spam. Separately, the administrator makes the \emph{disposition}: Accepted, Pending Clarification or Rejected with documented reason. Figure~\ref{fig:op-hitl} shows the decision boundary. A predicted Spam that was actually a clear complaint can be corrected to Valid and accepted; an incomplete original can remain Suspicious while external clarification permits administrative acceptance. No accepted case is automatically converted to a Valid training label, and no rejected case is automatically converted to Spam. A clarification request must be reassessed after a response or the applicable policy-governed response period; no invented universal deadline is imposed \cite{rottger2022,fleisig2024,franc2023}.
\begin{figure}[H]\centering\includegraphics[width=.96\textwidth,height=.83\textheight,keepaspectratio]{figures/hitl_proposed.pdf}\caption{Proposed Human-in-the-Loop complaint review with distinct text-label verification and administrative disposition. Only predicted Suspicious or Spam enter routine HITL; direct Valid predictions proceed separately.}\label{fig:op-hitl}\end{figure}

\subsection{Nine Boundary Scenarios for Consistent Labeling}
Table~\ref{tab:op-boundary} operationalizes the ambiguous cases found in the consolidated evidence review. The text label applies to the \emph{narrative version}, not the independently investigated incident or the administrative case outcome. These are annotation and workflow rules to be validated by reviewers, not performance findings from an authentic dataset \cite{rottger2022,fleisig2024,li2023}.
\begin{longtable}{p{.39\textwidth}p{.23\textwidth}p{.29\textwidth}}
\caption{Proposed boundary-case labeling and administrative handling rules}\label{tab:op-boundary}\\
\toprule Scenario & Verified narrative label & Administrative handling\\\midrule\endfirsthead
\toprule Scenario & Verified narrative label & Administrative handling\\\midrule\endhead
Predicted Suspicious, but original text is already clear. & Correct to Valid. & Accept and route to case management.\\
Original is vague; resident supplies complete revision. & Original Suspicious; revised Valid. & Accept revised narrative; preserve both versions.\\
Original vague; independent evidence supports proceeding. & Original remains Suspicious. & May accept if authorized procedure permits.\\
Predicted Spam, but original is an actionable complaint. & Correct to Valid. & Accept; do not automatically reject.\\
Predicted Spam, but original is complaint-related and incomplete. & Correct to Suspicious. & Clarify; accept if sufficient information follows.\\
Predicted Suspicious, but content is promotional non-complaint material. & Correct to Spam. & Reject after human confirmation and record reason.\\
Text is clear but allegation has not been substantiated. & Valid as text. & Investigate or document outcome; truth not inferred by SVM.\\
Clear text later has legitimate administrative grounds for closure or rejection. & Valid as text. & Record grounds without changing text label.\\
Repeated similar submissions from one or several residents. & Label each version by its text. & Check duplicates and corroboration independently.\\
\bottomrule\end{longtable}

\subsection{Verified Feedback, Data Quality and Controlled Model Retraining}
Figure~\ref{fig:op-retrain} begins with versioned human-reviewed label records; it does not repeat the administrative Accept/Reject workflow. Records first undergo quality, duplicate, provenance, privacy and split-family checks. They are then accumulated in a verified-feedback pool. A periodic review of count, verified-class distribution, diversity, repeated errors and performance degradation determines whether a candidate training run is justified; the run requires authorization. Retraining is \emph{not triggered by each click} on an administrative button. An approved run consolidates the existing eligible training material and the new verified examples, refits the preprocessing and vectorization components, trains and recalibrates a candidate SVM, and evaluates it on a fixed independent validation protocol. A candidate that fails release criteria is archived without replacing the live model; a candidate that passes can be separately approved and versioned for future predictions \cite{figueroa2012,ghose2024,sklearnpitfalls,nist}.

The consolidated research review proposes cumulative increments of \emph{100, 250 and 500} verified examples as learning-curve checkpoints, with 1,000 only if feasible. They are \emph{research planning choices}, not empirically validated minimums or automatic deployment triggers. The baseline plus those three candidates are four conceptual model configurations, but cross-validation and calibration can require many additional fits. All variants of one case and related synthetic templates must remain grouped to prevent evaluation leakage. A small independent audit of predicted Valid cases should prevent the feedback pool from consisting solely of Suspicious/Spam predictions. Existing synthetic-data Macro-F1 scores cannot demonstrate that retraining on authentic complaints improves generalization \cite{figueroa2012,beatty2018,citovsky2021,ghose2024,lakkaraju2017}.
\begin{figure}[H]\centering\includegraphics[width=.96\textwidth,height=.83\textheight,keepaspectratio]{figures/retrain_proposed.pdf}\caption{Proposed quality-checked human feedback and controlled batch retraining; training eligibility and deployment approval are separate decisions, and failure does not trigger infinite retries.}\label{fig:op-retrain}\end{figure}

\subsection{Accountable Complaint Case Management (Accepted Cases Only)}
The accepted-case workflow in Figure~\ref{fig:op-case} begins after the system accepts a predicted Valid narrative or an authorized administrator accepts a reviewed submission. The case receives a tracking identifier and a principal state of \emph{Unresolved}; designation of responsible staff, documented investigation and progress updates move the case to \emph{Ongoing}; independently checked completion allows it to become \emph{Resolved}. Records capture officer assignment, action timestamps, correspondence and supporting justification. \emph{Overdue}, \emph{Referred} and \emph{Reopened} are auditable events or flags, not additional SVM labels. If an action is unfinished, disputed or inadequately verified, the case returns to follow-up rather than being silently finalized.

This proposed administrative design is grounded in the distinction between public \emph{response obligations} and full incident resolution. Republic Act No.~6713, Section~5(a), addresses replies to public communications within fifteen working days; Republic Act No.~11032 and the relevant Citizen's Charter govern processing of covered services by complexity. The special 72-hour concrete-action standard associated with the national 8888 channel is \emph{not} a universal deadline for resolving every barangay incident. The application must store the locally applicable rule, permit reasoned referral/escalation and support resident disagreement. Davao City Reports provides a local precedent for tracking references, progress history and re-filing, not a claim that this proposed barangay implementation already meets its requirements \cite{ra6713,ra11032,eo8888,dcr}.
\begin{figure}[H]\centering\includegraphics[width=.97\textwidth,height=.87\textheight,keepaspectratio]{figures/case_proposed.pdf}\caption{Proposed accountable complaint case management: accepted-case tracking, action logging, policy-based follow-up, completion verification, resident feedback and justified reopening.}\label{fig:op-case}\end{figure}

\subsection{End-to-End Integrated Complaint Workflow}
The integrated workflow in Figure~\ref{fig:op-integrated} links the previously described modules into one end-to-end account of proposed system behavior. A complaint narrative passes through conservative NLP preprocessing, saved word/character TF--IDF transformers and a calibrated SVM. Predicted Valid narratives proceed to accepted-case management; only predicted Suspicious or Spam narratives enter routine HITL for text-label verification and a separate administrative decision. Reviewed but non-accepted submissions retain a documented clarification or rejection outcome. \emph{Accepted} is an administrative state, not a replacement for the narrative label, while \emph{Resolved} is a later case-management state, not proof that an originally Suspicious sentence was Valid.

The diagram intentionally separates live case handling from the offline learning loop. Case records store assignment, action history, response deadlines and resolution outcomes, whereas versioned human-verified narrative labels are quality-checked before they can enter the training-feedback pool. Training a new TF--IDF--SVM candidate requires a justified and authorized retraining decision; the deployed classifier changes only after successful validation and release approval. In particular, neither an Accept click nor case resolution automatically updates model weights or training labels. These are proposed control boundaries, not additional observed results from the recorded synthetic-data baseline \cite{ra10173,rottger2022,sklearnpitfalls,nist}.
\begin{figure}[H]\centering\includegraphics[width=.99\textwidth,height=.87\textheight,keepaspectratio]{figures/end_to_end_workflow.pdf}\caption{Proposed end-to-end integrated workflow connecting resident submission, NLP, TF--IDF, SVM, class-specific human review, accepted-case processing and approved offline retraining. Dashed arrows distinguish training-feedback and approved model updates from live complaint handling.}\label{fig:op-integrated}\end{figure}

\subsection{Software Component Integration and Data Boundaries}
After the end-to-end workflow, Figure~\ref{fig:op-system-integration} presents the complementary \emph{software component integration} and its access boundaries. The proposed resident portal and authorized administrator dashboard exchange data through a FastAPI application layer rather than calling the classifier or writing database records directly. The backend manages access checks, request validation, classification requests, tracking references, administrative actions and auditable case updates. Next.js/TypeScript, FastAPI/Python and PostgreSQL are the suggested application stack; the uploaded reproducibility evidence does not demonstrate deployment of these services.

The text-only inference service consumes the complaint narrative, applies the saved compatible preprocessing and TF--IDF transformations, and returns a saved calibrated SVM's proposed class probabilities. Predicted Valid follows the accepted-case service; predicted Suspicious or Spam enters administrator review. Reviewers may correct the \emph{text label} independently of Accept/Clarify/Reject decisions. Accepted cases and non-accepted submission outcomes share secured persistence but maintain different administrative states, while eligible verified narrative labels are stored as a \emph{distinct feedback dataset}. Later, an authorized offline candidate training run consolidates eligible training records, retrains the feature extractors and classifier, checks calibration and held-out validation, and releases a single versioned model bundle only after approval. The dashed release path denotes an approved future inference version, never incremental training during complaint submission. The proposed design excludes photographic inputs and location maps from model inference and does not claim that a completed case is necessarily a Valid textual training example \cite{sklearnpitfalls,sklearncalib,ra10173,nist}.
\begin{figure}[H]\centering
\includegraphics[width=.99\textwidth,height=.84\textheight,keepaspectratio]{figures/verified/system_architecture.pdf}
\caption{Proposed system integration architecture showing the resident and administrator interfaces, application API, text-classification service, HITL review, separate case and feedback records, and controlled offline model-release boundary. The component diagram specifies a target implementation, not a deployed or evaluated full-stack system.}\label{fig:op-system-integration}
\end{figure}

\subsection{Implementation and Evaluation Gates}
\begin{table}[H]\centering\small\caption{Evidence boundary and evaluation required for proposed extensions}\label{tab:op-evidence}\begin{tabularx}{\textwidth}{p{.29\textwidth}Xp{.24\textwidth}}\toprule
Component & What the supplied research currently supports & Required next evidence\\\midrule
24,000-row synthetic dataset & Fixed counts, partitions, metadata and researcher-generated texts; all review flags negative. & Human annotation and authentic external validation.\\
Word TF--IDF + Linear SVM & Recorded group-aware grid/CV and within-generator validation; test generalization remains limited. & Independent test beyond generator patterns.\\
Word + character combination & Reference reproduction and separate combined-feature development experiment. & Independent external comparison beyond the synthetic ceiling.\\
Probability calibration and HITL & Synthetic development calibration; HITL remains a proposed policy. & External probability reliability, reviewer agreement and Valid-sample audits.\\
Human-feedback retraining & Cumulative checkpoints and release-gate plan. & Verified feedback, fixed validation, approved candidate run.\\
Case management/accountability & Flowchart and local legal/interface precedents. & Implemented audit logs, deadlines, referral and resident tests.\\\bottomrule
\end{tabularx}\end{table}

\subsection{Detailed Full-Stack Implementation Plan}
The separately delivered roadmap defines the Next.js/TypeScript, FastAPI and PostgreSQL modules, REST schemas, versioned narrative/label/disposition/case tables, security controls, phased tests and acceptance gates. Figure~\ref{fig:op-system-integration} supplies the single software-integration overview; Figure~\ref{fig:verified-db} gives the complementary relational model. These are implementation plans, not validation of the existing SQLite/static-HTML prototype. Direct Valid routing records policy-based acceptance but never creates a human-verified label. Urgent human handling is independent of predicted class. The full database cardinalities and editable diagram sources accompany the roadmap.
The relational design links immutable narrative versions to separate predictions and verified labels; accepted dispositions alone create cases. Append-only action and audit histories preserve the distinction between machine outputs and administrative decisions.

\begin{figure}[H]\centering
\includegraphics[width=.96\textwidth,height=.70\textheight,keepaspectratio]{figures/verified/database_overview.pdf}
\caption{Proposed logical relational overview; full keys, constraints, indices and cardinalities are supplied with the roadmap.}\label{fig:verified-db}
\end{figure}


\section{Research Ethics, Reproducibility and Limitations}
The corpus contains researcher-generated narratives rather than private resident reports. Reproducible inputs include the category register, generator source, random seed, partition CSVs, run records, training scripts and model metrics. Ethical deployment requires privacy-conscious logging and retention, a clear review process, and transparent handling of urgent reports. Because generator-authored patterns remain correlated with class labels, the measured within-generator accuracy can exaggerate transferable performance. No prospective performance gain from Human-in-the-Loop retraining or reduction in office workload is claimed without an implemented and evaluated user study.

\section{Executed Rerun Environment and Separate Candidate Protocol}
The completed reproduction used Python 3.13.5, scikit-learn 1.8.0, NumPy 2.4.4, SciPy 1.17.1, pandas 3.0.2, Matplotlib 3.11.2 and joblib 1.5.3 on Windows AMD64. The workstation specifications earlier in this chapter are supplied historical specifications rather than independently measured hardware claims. Seed 20261026 was retained. The source SHA-256 was verified before and after execution, and original outputs were preserved before rerunning extracted copies. The local portable input is \path{data/Barangay_Complaint_Dataset_24000.csv}. Only \texttt{complaint\_text} enters the predictor; metadata remain excluded.

Historical reference text is passed directly to the saved scikit-learn vectorizer: default lowercasing/tokenization, no stopword removal, word n-grams, sublinear term frequency, smooth inverse document frequency and L2 normalization. The demo script's URL/punctuation-stripping helper was not retroactively inserted into the reference. The baseline uses unigrams/bigrams, \texttt{min\_df=2}, \texttt{max\_df=0.95}, and MultinomialNB $\alpha=1$. The chosen SVM uses unigrams with the same frequency/weighting settings and LinearSVC $C=0.1$. This compares complete pipelines with differing feature ranges.

After original reproduction, a single fixed combined-feature candidate applies Unicode NFKC, casefolding and whitespace normalization, without removing negations, Unicode accents, URLs or punctuation beforehand. Word TF--IDF uses $(1,2)$ n-grams, \texttt{min\_df=2}, \texttt{max\_df=0.95}; character-within-word TF--IDF uses $(3,5)$ n-grams, \texttt{min\_df=3} and a 45,000-feature cap. Both branches use sublinear weighting and default L2 normalization; their equally weighted sparse matrices are concatenated before LinearSVC $C=0.1$. No new candidate hyperparameter search was performed.

The calibrated candidate uses \texttt{CalibratedClassifierCV(method=sigmoid, ensemble=True)} with three precomputed \texttt{StratifiedGroupKFold} training-only splits. Each base feature pipeline is fitted inside its own fit fold, with its disjoint held-out group fold used for sigmoid fitting. Validation evaluates the ensemble; it never fits vocabulary, classifier or calibrator. The existing validation set is a development comparison, not new external confirmation. Candidate test evaluation was deliberately excluded. Log loss, multiclass Brier score (sum over classes), and ten equal-width confidence-bin expected calibration error describe development probability reliability. Report scores do not establish field-calibrated probabilities.

All final experimental figures are reproducibly generated from actual outputs at 360 DPI as PNG and also exported as vector PDF/SVG. CSV predictions and fold records remain machine-readable audit evidence, not the final evaluation presentation. The manuscript and standalone report use figures and academic tables with LaTeX captions and labels.
\par Macro-precision and macro-recall are $P_{\mathrm{macro}}=\frac{1}{3}\sum_c P_c$ and $R_{\mathrm{macro}}=\frac{1}{3}\sum_c R_c$. With no discordant pairs, the paired statistic is set to zero and the exact p-value to one.
